\paragraph{Identidad de la sección y composición de sus observables.}
\label{el-ampl-articulacion}
La identidad estructural del electrón se introduce mediante la sección
central y sus operaciones de transporte. La inversión de la coordenatización APP
selecciona la celda $(5,5)$ por su propiedad de punto fijo. Las dos hojas
aritméticas permiten estudiar las deformaciones que conservan la suma y
el residuo del producto. El resultado distingue las dos orientaciones
$(8,2)$ y $(2,8)$ y conserva la diferencia de una unidad de cociente.
Éste es el primer nivel de información: una celda, una involución, una
fibra conservada y un defecto entero mínimo. Ninguna evaluación de masa
es necesaria para enunciarlo o demostrarlo.

El paso a las proyecciones aritméticas introduce un segundo nivel. Las
fibras de suma y producto determinan esperanzas condicionales sobre el
mismo espacio finito. Su conmutador mide la dependencia respecto del
orden de esas dos proyecciones. La norma $\sqrt3/4$ resulta del espectro
de su emparejamiento; el bloque que alcanza esa norma conserva además
las relaciones cuadráticas demostradas a continuación. Por tanto, el
prefactor escalar y el álgebra local tienen una procedencia común, pero
contienen información diferente. La norma retiene la magnitud del
conmutador, mientras las matrices conservan su orientación y su ley de
composición.

La transferencia areal--volumétrica incorpora las coordenadas regionales
en un tercer nivel. El calendario de modos determina la matriz de
incidencia y su acción cuadrática. La marca de clausura produce la razón
$\pi/\varphi^2$; el intercambio de los dos argumentos proporciona la
orientación $\varphi/\pi^2$ utilizada en el carácter exponencial
electrónico. Las dos expresiones quedan relacionadas por una involución
explícita. Esa relación explica por qué una coordenada regional puede
aparecer en la ecuación electrónica con una potencia o en una posición
diferente de la que tenía en la primera lectura: la operación aplicada
es parte del dato matemático que se conserva.

El término cúbico de alfa pertenece al complemento regional de un retorno
marcado. La cuenta de sus posiciones produce $27-5=22$, mientras las
tres hojas de cada una de las cinco regiones producen $3\cdot5=15$.
Las demostraciones conservan el dominio de esos recuentos y la
orientación que determina sus signos. La fórmula basal compone,
por consiguiente, una norma de incompatibilidad, un carácter de
transferencia, un déficit de posiciones y una memoria torsional.
La escritura en una sola línea expresa la composición de esas
operaciones; su explicación exige reconstruirlas en ese orden.

La transición de la coordenada basal a la coordenada completa conserva
la misma sección central. El factor de retorno utiliza los registros
de la historia, y los dos lectores declarados recuperan el triple
$(80,54,6)$ sobre esa trayectoria. La multiplicación por
$R_{\mathrm{act}}^{80-54\alpha+6\Delta_4}$ modifica la evaluación,
pero mantiene identificados su origen y sus antecedentes. El valor
basal constituye una etapa intermedia de la composición; el valor
completo incluye también esa contribución del retorno. La comparación
física debe utilizar la etapa especificada por la fórmula que se evalúa.

La escala de acción interviene después con otra función. Una razón de
dos secciones de acción es adimensional y cancela su escala común.
La realización absoluta utiliza, en cambio, una base de acción, un
lector temporal y el mapa hacia energía o masa. Esta separación permite
conservar simultáneamente la construcción de la escala y la identidad
de la sección electrónica. El cuanto cronológico proporciona una
referencia de escala; la sección central determina el factor estructural
que se expresa respecto de ella. Las ecuaciones dimensionales de este
capítulo mantienen explícita esa composición.

La pertinencia del electrón en el argumento general procede de esta
convergencia de operaciones. La misma aritmética que produjo las
regiones determina ahora una celda central, sus deformaciones y sus
proyecciones. La constante de estructura fina, construida antes, actúa
en esa composición como coordenada de clausura. La información de
retorno determina la evaluación completa, y la realización dimensional
permite el contraste posterior. Explicar el electrón desde la estructura
consiste en conservar esta sucesión de objetos, mapas y valores en cada
paso de la demostración.
